\subsection{实验现象与问题分析}

本次实验围绕 ARP 缓存污染和中间人攻击展开。在实验过程中，
不同攻击方式表现出了不同的网络现象。

首先，在 ARP Request 和 ARP Reply 缓存污染实验中，
当 Victim 的 ARP 缓存被成功污染后，
其发往目标主机的数据帧在 IP 层的目的地址并未发生改变，
但数据链路层的目的 MAC 地址已经被替换为 Attacker 的 MAC 地址。
这说明 ARP 攻击本质上并不改变网络层的通信目标，
而是通过破坏 IP 地址与 MAC 地址之间的映射关系，
改变以太网帧在局域网中的实际转发方向。

其次，在仅完成 ARP 污染、但没有开启 Attacker 的 IP 转发功能时，
Victim 发出的数据能够到达 Attacker，
但 Attacker 不会继续将其转发至真实目标，
因此出现 ping 超时或大量丢包的现象。
这说明 ARP 缓存污染本身只能实现流量重定向，
若要形成持续且不易被用户察觉的中间人攻击，
还需要攻击者进一步承担流量转发功能。

在 ARP Reply 实验过程中，还观察到了真实 ARP 响应与伪造 ARP 响应之间的竞争。
由于正常主机仍可能发送正确的 ARP 信息，
被污染的邻居项可能重新恢复。
因此，攻击者通常需要持续发送伪造的 ARP Reply，
以维持错误的 IP--MAC 映射关系。

最后，在中间人攻击实验中，
开启 Attacker 的 IP 转发功能后，
Victim 与 Web Server 之间的通信仍然保持正常。
这说明用户仅根据网络连接是否正常，
往往难以直接发现自己的流量已经经过第三方主机，
体现了 ARP 中间人攻击具有一定隐蔽性。

\subsection{安全防护与改进思考}

本实验表明，ARP 协议本身缺乏对报文来源真实性的认证，
攻击者能够通过伪造 ARP Request 或 ARP Reply
改变局域网主机的 IP--MAC 映射关系。
因此，在实际网络环境中可以采用多种机制降低 ARP 攻击风险。

首先，可以在交换机上部署 Dynamic ARP Inspection（DAI），
并结合 DHCP Snooping 建立可信的 IP--MAC--端口绑定关系，
对不符合绑定关系的 ARP 报文进行过滤。
对于网络结构固定、设备数量较少的重要主机，
也可以采用静态 ARP 绑定，
减少错误 ARP 映射被写入缓存的可能性。

其次，可以对局域网中的 ARP 映射变化进行持续监测。
如果同一个 IP 地址在短时间内频繁对应不同 MAC 地址，
或者网关 MAC 地址异常变化，
则可能说明网络中存在 ARP 欺骗行为。

对于上层通信，本实验中的 HTTP Digest 认证虽然避免了密码明文直接传输，
但攻击者一旦获得完整的认证参数，
仍可能针对弱口令实施离线猜测。
因此，在实际系统中应优先采用 HTTPS/TLS 对通信链路进行加密，
同时使用足够复杂的高强度密码，
必要时结合多因素认证机制，
从而降低中间人窃听和离线口令猜测带来的风险。

通过本次实验可以看出，
网络安全不能仅依赖某一层协议提供保护，
而应同时从二层网络访问控制、传输链路加密、
身份认证和异常检测等多个层面建立纵深防御机制。
